% What happens once and what happens again. Referenced from chapter 14. Bare
% tikzpicture; the figure environment, caption and label live at the call site.
%
% The point of the picture is the asymmetry between the two rows. The trigger is
% a single request at the left edge and never recurs; every fetch under it is
% paid again for each message, so the cost of a subscription is set by how much
% of the payload the trigger's own subgraph cannot answer.
\begin{tikzpicture}[
    font=\small,
    dot/.style={circle, draw, fill=black, inner sep=1.1pt},
    hollow/.style={circle, draw, fill=white, inner sep=1.1pt},
    flow/.style={-{Stealth[length=2mm]}},
    note/.style={font=\scriptsize, gray, align=center, inner sep=1pt},
    lbl/.style={font=\scriptsize, align=right, inner sep=2pt}
  ]

  \draw[flow] (1.4,0) -- (10.4,0);
  \node[note, anchor=west] at (10.5,0) {time};

  % -- the trigger, once ------------------------------------------------------
  \node[lbl, anchor=east] at (1.3,0.9) {trigger};
  \node[dot] (trig) at (2.0,0.9) {};
  \node[note, anchor=south west, align=left] at (2.1,1.05)
    {one \code{subscribe} to \code{reviews},\\then nothing};

  % -- the events -------------------------------------------------------------
  \node[lbl, anchor=east] at (1.3,0) {events};
  \foreach \x in {3.4,4.8,6.2,7.6,9.0}{
    \node[hollow] at (\x,0) {};
  }

  % -- the fetch each one pays for --------------------------------------------
  \node[lbl, anchor=east] at (1.3,-1.1) {\code{accounts}};
  \foreach \x in {3.4,4.8,6.2,7.6,9.0}{
    \node[dot] (f\x) at (\x,-1.1) {};
    \draw[flow] (\x,-0.15) -- (\x,-0.95);
  }

  \node[note, anchor=north, align=center] at (6.2,-1.4)
    {one \code{\_entities} call per message, for the one field\\the subgraph
     holding the connection cannot answer};

  \draw[dashed, gray] (2.0,0.7) -- (2.0,-1.6);
  \node[note, anchor=north east, align=right] at (1.95,-1.6) {measured once};

\end{tikzpicture}
